\documentclass[UTF8,a4paper,11pt]{ctexart}

\usepackage{amsmath,amssymb}
\usepackage{booktabs}
\usepackage{longtable}
\usepackage{geometry}
\usepackage{hyperref}
\usepackage{enumitem}
\usepackage{xcolor}
\usepackage{array}

\geometry{left=2.4cm,right=2.4cm,top=2.4cm,bottom=2.4cm}
\hypersetup{colorlinks=true,linkcolor=blue!60!black,urlcolor=blue!60!black,citecolor=blue!60!black}

\newcommand{\BER}{\operatorname{BER}}
\newcommand{\sBER}{\operatorname{sBER}}
\newcommand{\argmax}{\operatorname*{arg\,max}}
\newcommand{\argmin}{\operatorname*{arg\,min}}

\title{Global Guarded Vector Hue 位置指纹通信实验流程说明}
\author{Position Fingerprint Experiment}
\date{2026年5月14日}

\begin{document}
\maketitle

\begin{abstract}
本文档根据 \texttt{vector\_hue/global\_guarded\_core.py}、\texttt{virtual\_stream\_core.py} 和当前实验输出文件整理完整实验流程。实验目标不是给设备人工分配正交码，而是从不同空间位置的实测响应中提取位置指纹；发送端把多路 bit 组合映射为四通道 vector hue；合法设备使用自己的位置指纹正确恢复自己的信息；非法设备即使用自己的位置指纹尝试解码，也只能得到接近随机的结果。文档包含基于真实数据的完整例子：合法位置为 $(2,5,26)$，候选编号为 267，真实 Probe 集合为 $[25,75,100,115,180,215,225,250,270,295,315,325,335,340,355]$。
\end{abstract}

\section{一句话概括}

当前实验可以概括为：\textbf{从每个位置的实测 Probe 响应中提取本地地址码，发送端用 vector hue mapping 同时承载多路信息，合法位置用自己的地址码做相关解码，非法位置由于指纹不同而无法稳定恢复任一路合法信息。}

与 PDF 1 中的两位置最小 PoC 相比，当前代码已经扩展为更完整的安全实验：
\begin{enumerate}[label=(\arabic*)]
  \item 支持 $k$ 个真实合法位置，本例 $k=3$；
  \item 通过虚拟流把有效维度扩展到 \texttt{target\_effective\_k=8}；
  \item 每个符号组合不再映射到单个 hue，而是映射到四个 chip 各自的 hue，即 vector hue；
  \item 不只追求合法端正确，还显式搜索并修复非法端最容易泄露的 weak routes；
  \item 最后用线性规划混合多个候选 schedule，使最坏非法路由也尽量接近随机猜测。
\end{enumerate}

\section{实验对象和基本名词}

\subsection{位置、Probe 和 chip}

实验数据存放在 \texttt{data/15pro/yellow\_shuffled}。每个 CSV 文件对应一个空间位置，例如位置 2、位置 5、位置 26。每一行对应一个 Probe hue，例如 25、75、100；每一行有 4 列数值，表示四个显色片或四个 chip 的响应。因此，对位置 $p$ 和 $m$ 个 Probe，可以得到响应矩阵
\[
Y_p\in\mathbb{R}^{m\times 4}.
\]
本例中 $m=15$，四个 chip 分别记为 chip1 到 chip4。

\subsection{合法设备和非法设备}

一次实验选出一组合法位置
\[
\mathcal{L}=\{\ell_1,\ell_2,\ldots,\ell_k\}.
\]
这些位置是目标接收者。其他位置视为非法位置
\[
\mathcal{I}=\mathcal{P}\setminus\mathcal{L}.
\]
本例中合法位置为
\[
\mathcal{L}=\{2,5,26\}.
\]
位置 25 是非法位置之一，后文用它举例说明非法解码失败。

\subsection{真实流和虚拟流}

\texttt{real\_k=3} 表示真实需要服务的合法位置有 3 个。代码还设置
\[
\texttt{target\_effective\_k}=8,
\]
所以会补充
\[
\texttt{virtual\_stream\_count}=8-3=5
\]
个虚拟流。虚拟流不对应真实用户，它们的作用是增加符号组合空间，使非法端更难把自己的输出稳定对应到某一路真实 bit。

\section{位置指纹如何从真实数据中长出来}

\subsection{线性去趋势}

对每个位置 $p$，输入 Probe 序列 $x=[q_1,\ldots,q_m]$ 和四通道响应矩阵 $Y_p$。代码在 \texttt{extract\_fingerprint} 中对每个 chip 单独拟合线性趋势：
\[
Y_p[:,d]\approx a_{p,d}x+b_{p,d}.
\]
然后得到残差矩阵
\[
R_p=Y_p-T_p.
\]
直观理解：线性趋势表示 hue 变大时响应整体变大的公共部分；残差表示这个位置独有的非线性起伏，这部分更像位置指纹。

\subsection{SVD 提取读出方向}

对残差矩阵做奇异值分解：
\[
R_p=U\Sigma V^\top.
\]
取第一右奇异向量作为读出方向：
\[
w_p=V_1.
\]
然后把每个 Probe 的四通道残差压缩成一维值：
\[
z_p=R_pw_p.
\]
最后取符号得到本地地址码：
\[
c_p[t]=\begin{cases}
+1,&z_p[t]\ge 0,\\
-1,&z_p[t]<0.
\end{cases}
\]

\subsection{真实数据例子：三个合法位置的指纹}

本例使用候选 267 的 15 个 Probe：
\[
[25,75,100,115,180,215,225,250,270,295,315,325,335,340,355].
\]
从真实 CSV 中读取响应并计算后，得到如下位置指纹。

为避免长向量在表格中挤压，下面按位置分别列出。每个位置的 $w_p$ 是四个 chip 的读出方向，$c_p$ 是 15 个 Probe 对应的地址码，$z_p$ 是残差投影值。

\paragraph{位置 2}
\[
w_2=[0.1300,\ 0.5312,\ -0.7659,\ -0.3380]
\]
\[
c_2=[+,-,+,-,+,-,-,+,-,+,+,+,+,+,-]
\]
\begin{align*}
z_2\approx[&213.356,\ -163.890,\ 127.067,\ -106.743,\ 28.668,\\
&-200.094,\ -240.826,\ 100.507,\ -51.133,\ 110.077,\\
&165.079,\ 67.521,\ 21.722,\ 75.131,\ -146.441].
\end{align*}

\paragraph{位置 5}
\[
w_5=[-0.6661,\ 0.2772,\ -0.5661,\ 0.3987]
\]
\[
c_5=[+,+,-,-,+,+,-,-,+,-,+,-,-,+,+]
\]
\begin{align*}
z_5\approx[&243.978,\ 17.414,\ -178.100,\ -115.487,\ 3.541,\\
&64.498,\ -91.178,\ -161.965,\ 115.029,\ -45.694,\\
&17.215,\ -2.594,\ -100.079,\ 87.252,\ 146.171].
\end{align*}

\paragraph{位置 26}
\[
w_{26}=[-0.2867,\ -0.0525,\ -0.8452,\ -0.4481]
\]
\[
c_{26}=[+,-,+,+,-,+,-,+,-,+,-,+,-,+,+]
\]
\begin{align*}
z_{26}\approx[&8.190,\ -76.789,\ 70.660,\ 83.474,\ -224.310,\\
&118.463,\ -4.068,\ 212.320,\ -198.156,\ 58.092,\\
&-221.409,\ 162.834,\ -100.011,\ 65.344,\ 45.366].
\end{align*}

这些码不是人工设计的 Hadamard 码，而是由三个位置的真实响应自动产生的。

\section{发送端如何把 bit 编成 vector hue}

\subsection{从 bit 到符号组合}

每个发送块包含 8 路符号：前三路是真实用户，后五路是虚拟流。记为
\[
\mathbf{b}=(b_1,b_2,\ldots,b_8),\quad b_j\in\{-1,+1\}.
\]
第 $t$ 个 Probe 时隙的符号组合为
\[
\mathbf{s}[t]=(b_1c_1[t],b_2c_2[t],\ldots,b_8c_8[t]).
\]
其中虚拟流复用真实位置的码，代码中由 \texttt{full\_codes\_for\_virtual} 完成。

\subsection{什么是 vector hue}

普通 hue mapping 是
\[
M:\{-1,+1\}^k\rightarrow q,
\]
即一个符号组合对应一个 hue。当前实验使用 vector hue mapping：
\[
M:\{-1,+1\}^{8}\rightarrow(q_1,q_2,q_3,q_4),
\]
四个分量分别送给四个 chip。也就是说，在同一个时隙内，chip1 可以使用 hue 100，chip2 可以使用 hue 215，chip3 可以使用 hue 75，chip4 可以使用 hue 295。

\subsection{Hue mapping 如何计算}

对某个符号组合 $\mathbf{a}=(a_1,\ldots,a_8)$，代码先看真实合法位置的前三个符号。对每个 chip $d$ 和候选 Probe $q$，计算得分
\[
\operatorname{score}_d(q)=\sum_{i=1}^{3} a_i R_{\ell_i}[q,d]w_{\ell_i}[d].
\]
得分越高，说明这个 chip 选择该 hue 后越能推动合法位置朝目标符号方向变化。每个 chip 取高分候选，再组合成四维 vector hue。代码位置为 \texttt{\_build\_vector\_candidate\_dict}。

\subsection{真实计算例子：符号 $[-1,-1,-1,-1,-1,-1,-1,-1]$}

候选 267 对该符号选出的 vector hue 为
\[
M([-1,-1,-1,-1,-1,-1,-1,-1])=(100,215,75,295).
\]
它不是手工指定的，而是按上面的得分公式计算得到。四个 chip 的真实得分如下。

\begin{center}
\begin{tabular}{ccccc}
\toprule
chip & 选择 hue & 位置2贡献 & 位置5贡献 & 位置26贡献 / 合计\\
\midrule
1 & 100 & -7.3962 & 106.0635 & 46.0768 / 144.7442\\
2 & 215 & 88.2653 & -3.8111 & 9.6599 / 94.1141\\
3 & 75 & 113.7853 & 0.3549 & 73.1553 / 187.2955\\
4 & 295 & 52.1155 & 34.6001 & -14.8979 / 71.8178\\
\bottomrule
\end{tabular}
\end{center}

四个 chip 合计得分约为
\[
144.7442+94.1141+187.2955+71.8178=497.9716.
\]
这说明该 vector hue 对目标符号方向总体是有利的。

\section{完整发送和合法解码算例}

\subsection{一个真实发送块}

从代码生成的第一个 block 为
\[
\mathbf{b}=(-1,-1,-1,-1,-1,-1,-1,-1),
\]
对应二进制为
\[
[0,0,0,0,0,0,0,0].
\]
前三路真实用户的目标 bit 都是 0，也就是 $(-1,-1,-1)$。

前五个时隙的符号组合和发送 vector hue 如下。

\begin{center}
\begin{tabular}{c>{\raggedright\arraybackslash}p{6.8cm}>{\raggedright\arraybackslash}p{4.2cm}}
\toprule
时隙 & 符号组合 $\mathbf{s}[t]$ & 发送 vector hue\\
\midrule
1 & $[-1,-1,-1,-1,-1,-1,-1,-1]$ & $[100,215,75,295]$\\
2 & $[1,-1,1,1,-1,1,1,-1]$ & $[225,25,250,250]$\\
3 & $[-1,1,-1,-1,1,-1,-1,1]$ & $[355,325,315,340]$\\
4 & $[1,1,-1,1,1,-1,1,1]$ & $[25,25,315,340]$\\
5 & $[-1,-1,1,-1,-1,1,-1,-1]$ & $[325,215,115,225]$\\
\bottomrule
\end{tabular}
\end{center}

注意：所有合法设备和非法设备接收到的是同一串公共 vector hue 序列。差别只在于每个位置的实测响应矩阵不同，所以同一发送序列在不同位置会变成不同观测。

\subsection{合法设备的本地解码公式}

位置 $p$ 收到一个 block 后，先根据自己的真实响应矩阵查表，得到本地观测
\[
Y_{p,\mathrm{obs}}\in\mathbb{R}^{15\times4}.
\]
在线解码时并不知道发送端真实 bit，因此不能减去已知趋势，只做块内中心化：
\[
\widetilde{Y}_{p}=Y_{p,\mathrm{obs}}-\bar{Y}_{p,\mathrm{obs}}.
\]
然后投影：
\[
u_p=\widetilde{Y}_pw_p.
\]
最后与本地地址码相关：
\[
\gamma_p=c_p^\top u_p.
\]
判决规则为
\[
\hat b_p=\begin{cases}
+1,&\gamma_p>0,\\
-1,&\gamma_p\le 0.
\end{cases}
\]

\subsection{真实合法解码结果}

对上面的 block，三个合法位置的真实计算结果如下。

\begin{center}
\begin{tabular}{ccccc}
\toprule
合法位置 & 目标符号 & $u_p$ 前五项 & $\gamma_p$ & 判决\\
\midrule
2 & $-1$ & $[-251.460,149.749,39.678,145.117,-232.389]$ & $-1771.5279$ & $-1$ 正确\\
5 & $-1$ & $[-102.894,-165.288,164.176,188.337,-173.011]$ & $-2050.0303$ & $-1$ 正确\\
26 & $-1$ & $[-53.007,275.969,-188.897,-149.776,134.064]$ & $-2310.4931$ & $-1$ 正确\\
\bottomrule
\end{tabular}
\end{center}

因为三个 $\gamma_p$ 都小于 0，三个合法设备都判为 $-1$，与发送目标完全一致。

\section{非法设备为什么解码失败}

\subsection{非法设备也可以尝试同样算法}

非法位置 25 也有自己的响应矩阵 $Y_{25}$、读出方向 $w_{25}$ 和地址码 $c_{25}$。它可以把同一串 vector hue 当作自己的观测来解码：
\[
\gamma_{25}=c_{25}^\top (Y_{25,\mathrm{obs}}-\bar{Y}_{25,\mathrm{obs}})w_{25}.
\]
但问题是：发送端的 mapping 是为了合法位置 $(2,5,26)$ 的指纹结构搜索出来的，不是为位置 25 搜索出来的。因此位置 25 的输出不会稳定等于任意一路合法 bit。

\subsection{单个 block 不足以说明安全，BER 才说明问题}

在第一个 block 中，非法位置 25 的计算为
\[
u_{25}\text{ 前五项}\approx[58.240,44.477,-1.939,0.233,-60.098],
\]
\[
\gamma_{25}\approx -460.6195,
\]
所以它也判为 0。这个 block 正好和合法 bit 一样，不能说明泄露，也不能说明安全。安全性要看 1000 个随机 block 中，它和每一路真实 bit 的误码率。

对候选 267，在 1000 个真实随机 block 上，非法位置 25 的结果为：
\begin{center}
\begin{tabular}{ccc}
\toprule
非法路由 & raw BER & corrected BER\\
\midrule
25 $\rightarrow$ 2 & 0.424 & 0.424\\
25 $\rightarrow$ 5 & 0.449 & 0.449\\
25 $\rightarrow$ 26 & 0.447 & 0.447\\
\bottomrule
\end{tabular}
\end{center}

这些值都接近 0.5，表示非法位置 25 对三路真实信息基本接近随机猜测。

\section{Global Guarded 搜索流程}

\subsection{第一阶段：baseline 候选}

\texttt{run\_global\_guarded\_experiment} 先随机抽样合法位置组合。对每个组合，生成多个 Probe 子集。每个 Probe 子集加载合法位置模型，再生成多个 hue mapping 候选。每个候选都必须满足：
\[
\BER_{\mathrm{authorized}}=0.
\]
也就是合法端在 1000 个 block 上不能出错。只有合法端完全正确的候选才进入安全评价。

\subsection{第二阶段：找 anchor 和 weak routes}

baseline 候选中，代码选出最好的 anchor：
\[
\text{anchor}=\argmax(\min \sBER_{\mathrm{route}},\ \operatorname{avg}\sBER_{\mathrm{route}}).
\]
然后找出该 anchor 最薄弱的若干非法路由，称为 weak routes。对本例，baseline 阶段记录为：
\begin{center}
\begin{tabular}{ll}
\toprule
字段 & 真实值\\
\midrule
合法位置 & $(2,5,26)$\\
baseline anchor candidate & 34\\
anchor 最小安全 BER & 0.26500000\\
anchor 最坏路由 & 20 $\rightarrow$ 5\\
weak routes & 20$\rightarrow$5, 24$\rightarrow$5, 9$\rightarrow$5, 22$\rightarrow$5, 22$\rightarrow$2, 28$\rightarrow$2, 27$\rightarrow$2, 7$\rightarrow$5\\
\bottomrule
\end{tabular}
\end{center}

这些 weak routes 告诉程序：当前最危险的是哪些非法位置试图恢复哪些合法信息。

\subsection{第三阶段：targeted repair}

代码随后进行 \texttt{anchor\_worst\_targeted} 搜索。它继续生成更多 Probe 子集和 mapping 候选，重点寻找能改善 weak routes 的候选。候选 267 就来自该阶段，来源字段为 \texttt{anchor\_worst\_targeted}，Probe 子集编号为 27。

\subsection{第四阶段：选择 20 个候选}

\texttt{select\_global\_guarded\_candidates} 不只选全局最强候选，而是组合四类候选：
\begin{enumerate}[label=(\arabic*)]
  \item anchor 本身，保证不低于一个已知强基准；
  \item 全局表现最好的候选；
  \item 对 weak routes 平均改善最大的候选；
  \item 对单条 weak route 有救援能力的候选；
  \item worst route 分布不同的候选，用来增加多样性。
\end{enumerate}
这样做的原因是：单个候选可能某些路由很强、某些路由很弱；多个候选轮换使用后，非法设备很难长期稳定利用同一个弱点。

\subsection{第五阶段：线性规划求 usage ratio}

假设最终选出 $n=20$ 个候选，每个候选在每条非法路由上有一个 raw BER。把这些值组成矩阵
\[
A\in\mathbb{R}^{R\times20}.
\]
代码用线性规划寻找使用比例
\[
\mathbf{r}=(r_1,\ldots,r_{20}),\quad r_j\ge 0,\quad \sum_j r_j=1,
\]
使混合后的最坏安全 BER 尽量大。混合 raw BER 为
\[
\mathbf{e}=A\mathbf{r}.
\]
安全 BER 使用 corrected BER：
\[
\sBER=\min(\BER,1-\BER).
\]
这是因为如果非法端 BER 等于 1，它只要取反就能完全恢复，所以 BER=1 和 BER=0 一样危险。

\section{本例最终结果}

本例来自 \texttt{vector\_hue/k3/global\_guarded/results\_summary.csv} 第一个样本，真实结果如下。

\begin{center}
\begin{tabular}{ll}
\toprule
字段 & 真实值\\
\midrule
real k & 3\\
effective k & 8\\
virtual stream count & 5\\
合法位置组合 & $(2,5,26)$\\
选中候选数 & 20\\
共同非法路由数 & 48\\
合法端最大 BER & 0.00000000\\
三个合法位置 BER & $[0.00000000,0.00000000,0.00000000]$\\
混合后最小安全 BER & 0.43411580\\
混合后平均安全 BER & 0.46757386\\
最终最坏路由 & 25 $\rightarrow$ 26\\
优化器 & linprog\\
anchor candidate & 267\\
anchor source & anchor\_worst\_targeted\\
anchor 最小安全 BER & 0.32300000\\
security gain over anchor & 0.11111580\\
\bottomrule
\end{tabular}
\end{center}

最终 schedule 长度为 1000，20 个候选的使用次数为：
\[
[240,156,0,0,274,0,0,35,84,0,0,0,0,73,0,13,0,125,0,0].
\]
这表示实际发送时不是永远使用一个 mapping，而是在 1000 个 block 中按上述次数轮换候选。这样可以把最坏非法路由从 anchor 单独使用时的 0.323 提高到全局混合后的 0.43411580，更接近理想随机猜测的 0.5。

\section{完整流程总结}

\begin{enumerate}[label=(\arabic*)]
  \item 读取所有位置的真实 CSV 响应数据。
  \item 随机抽样合法位置组合，例如 $(2,5,26)$。
  \item 为该组合生成多个 15-Probe 子集。
  \item 对每个合法位置、每个 Probe 子集做线性去趋势和 SVD，得到 $w_p,z_p,c_p$。
  \item 生成真实 bit 和虚拟 bit，形成 8 维符号组合。
  \item 根据残差、读出方向和目标符号，为每个符号组合搜索四通道 vector hue mapping。
  \item 用真实响应矩阵模拟合法端接收；合法端必须 1000 个 block 全部正确，BER 为 0。
  \item 用所有非法位置的响应矩阵模拟攻击；计算每个非法位置到每路合法 bit 的 raw BER 和 corrected BER。
  \item 从 baseline 候选中找 anchor 和 weak routes。
  \item 对 weak routes 做 targeted repair，继续搜索能修复弱路由的候选。
  \item 从所有合法候选中选出 20 个互补候选。
  \item 用线性规划求 usage ratio，使混合后的最坏非法安全 BER 最大。
  \item 输出 \texttt{results\_summary.csv}、\texttt{results\_selected.csv} 和 \texttt{weak\_routes.csv}。
\end{enumerate}

\section{结论}

本实验已经从 PDF 1 的两位置最小闭环扩展为更强的安全验证流程。真实数据样本 $(2,5,26)$ 表明：合法端三路 BER 均为 0；非法端最坏路由经过全局候选混合后 corrected BER 达到 0.43411580，已经接近随机猜测。关键点在于，地址码来自位置实测响应，vector hue mapping 由残差和 SVD 指纹计算得到，global guarded 机制再针对最危险非法路由做补强。

\end{document}
